% ============================================================
%  SCRIPT CHI TIẾT — AGENTIC BI (GROQ + LLAMA 3.3-70B)
%  Olist Data Lakehouse Platform
%  Compile: pdflatex -interaction=nonstopmode script_agentic_bi_groq.tex
% ============================================================
\documentclass[12pt, a4paper]{article}

% ── Encoding & Language ──────────────────────────────────────
\usepackage[utf8]{inputenc}
\usepackage[T5]{fontenc}
\usepackage[vietnamese]{babel}

% ── Layout ───────────────────────────────────────────────────
\usepackage[top=2.2cm, bottom=2.2cm, left=2.5cm, right=2.5cm]{geometry}
\usepackage{parskip}
\setlength{\parindent}{0pt}

% ── Typography ───────────────────────────────────────────────
\usepackage{lmodern}
\usepackage{microtype}
\usepackage{setspace}
\onehalfspacing

% ── Color ────────────────────────────────────────────────────
\usepackage[dvipsnames, table, x11names]{xcolor}
\definecolor{bronze}{RGB}{180, 103, 40}
\definecolor{silver}{RGB}{130, 140, 150}
\definecolor{gold}{RGB}{180, 145, 20}
\definecolor{groqgreen}{RGB}{0, 174, 114}
\definecolor{llamapurple}{RGB}{124, 58, 237}
\definecolor{demobox}{RGB}{15, 23, 42}
\definecolor{scriptbg}{RGB}{241, 245, 249}
\definecolor{titlebg}{RGB}{10, 15, 35}
\definecolor{codebg}{RGB}{30, 30, 46}
\definecolor{p4color}{RGB}{220, 38, 38}
\definecolor{intentblue}{RGB}{37, 99, 235}
\definecolor{sqlorange}{RGB}{234, 88, 12}
\definecolor{duckdbgold}{RGB}{200, 160, 0}
\definecolor{insightgreen}{RGB}{5, 150, 105}

% ── Boxes ────────────────────────────────────────────────────
\usepackage[most]{tcolorbox}
\tcbuselibrary{breakable, skins, listings}

% Speaker box
\newenvironment{speakerbox}[2]{%
  \begin{tcolorbox}[
    enhanced, breakable,
    colback=#1!6!white, colframe=#1,
    colbacktitle=#1, coltitle=white,
    fonttitle=\bfseries\large,
    title={#2}, arc=4pt, boxrule=1.5pt,
    left=8pt, right=8pt, top=6pt, bottom=6pt,
  ]
}{\end{tcolorbox}}

% Demo action box
\newtcolorbox{demobox}[1][]{
  enhanced, breakable,
  colback=demobox, colframe=cyan!60, coltext=white,
  fonttitle=\bfseries\small\color{cyan!80},
  title={DEMO — Thao tác màn hình},
  arc=4pt, boxrule=1pt, left=8pt, right=8pt, #1
}

% Script box
\newtcolorbox{scriptbox}{
  enhanced, breakable,
  colback=scriptbg, colframe=gray!40,
  arc=3pt, boxrule=0.8pt,
  left=12pt, right=8pt, leftrule=4pt,
}

% Code box (dark)
\newtcolorbox{codebox}[1][]{
  enhanced, breakable,
  colback=codebg, colframe=groqgreen!60, coltext=white,
  fonttitle=\bfseries\small\color{groqgreen},
  arc=4pt, boxrule=1pt, left=8pt, right=8pt,
  #1
}

% Phase box
\newtcolorbox{phasebox}[2]{
  enhanced,
  colback=#1!8!white, colframe=#1,
  colbacktitle=#1, coltitle=white,
  fonttitle=\bfseries\normalsize,
  title={#2}, arc=5pt, boxrule=1.2pt,
  left=8pt, right=8pt,
}

% Timing badge
\newcommand{\timing}[1]{%
  \hfill{\small\bfseries\color{gray}[#1]}%
}

% Phase label
\newcommand{\phaselabel}[2]{%
  \noindent\textcolor{#1}{\rule{4pt}{12pt}}\hspace{4pt}%
  \textbf{\textcolor{#1}{#2}}%
}

% ── Graphics & TikZ ──────────────────────────────────────────
\usepackage{tikz}
\usetikzlibrary{
  arrows.meta, positioning, fit, shapes.geometric,
  shapes.multipart, calc, backgrounds, shadows,
  decorations.pathreplacing, matrix, chains
}

% ── Tables ───────────────────────────────────────────────────
\usepackage{booktabs}
\usepackage{tabularx}
\usepackage{multirow}
\usepackage{array}

% ── Code listings ─────────────────────────────────────────────
\usepackage{listings}
\lstset{
  basicstyle=\ttfamily\footnotesize\color{white},
  keywordstyle=\color{groqgreen!90},
  stringstyle=\color{yellow!80},
  commentstyle=\color{gray!60},
  backgroundcolor=\color{codebg},
  frame=none, xleftmargin=8pt, xrightmargin=4pt,
  breaklines=true, breakatwhitespace=true,
  showstringspaces=false,
}

% ── Misc ─────────────────────────────────────────────────────
\usepackage{enumitem}
\usepackage{fontawesome5}
\usepackage{hyperref}
\hypersetup{colorlinks=true, linkcolor=groqgreen, urlcolor=groqgreen}

% ── Header / Footer ──────────────────────────────────────────
\usepackage{fancyhdr}
\pagestyle{fancy}
\fancyhf{}
\fancyhead[L]{\small\color{gray}Agentic BI — Groq + LLaMA 3.3-70B}
\fancyhead[R]{\small\color{gray}Olist Data Lakehouse}
\fancyfoot[C]{\small\color{gray}\thepage}
\renewcommand{\headrulewidth}{0.4pt}

% ============================================================
\begin{document}
% ============================================================

% ── Title Page ───────────────────────────────────────────────
\begin{titlepage}
\begin{tikzpicture}[remember picture, overlay]
  \fill[titlebg] (current page.north west) rectangle
    ([yshift=-8cm]current page.north east);
  % Groq accent stripe
  \fill[groqgreen!80] ([yshift=-8cm]current page.north west)
    rectangle ([yshift=-8.3cm]current page.north east);

  \node[anchor=north, text=white, font=\Huge\bfseries,
        yshift=-1.8cm] at (current page.north)
    {Agentic BI};
  \node[anchor=north, text=groqgreen!90, font=\Large\bfseries,
        yshift=-3.2cm] at (current page.north)
    {Groq API $\cdot$ LLaMA 3.3-70B Versatile};
  \node[anchor=north, text=white!70, font=\large,
        yshift=-4.5cm] at (current page.north)
    {Natural Language $\rightarrow$ SQL $\rightarrow$ Insight};
  \node[anchor=north, text=cyan!60, font=\normalsize,
        yshift=-5.7cm] at (current page.north)
    {Text-based SQL Pipeline $\cdot$ DuckDB $\cdot$ Gold Iceberg $\cdot$ Streamlit};
\end{tikzpicture}

\vspace{9cm}

\begin{tcolorbox}[
  colback=groqgreen!8, colframe=groqgreen!60,
  arc=6pt, boxrule=1.2pt
]
\renewcommand{\arraystretch}{1.6}
\begin{tabularx}{\textwidth}{lX}
\toprule
\textbf{Thành phần} & \textbf{Chi tiết} \\
\midrule
LLM Provider    & \textbf{Groq API} — LPU inference, sub-second latency \\
Model           & \texttt{llama-3.3-70b-versatile} \\
Mode            & Text-based SQL (JSON output, self-correction retries) \\
Query Engine    & DuckDB in-memory — đọc trực tiếp từ Gold Iceberg tables \\
Frontend        & Streamlit — chat UI + Plotly chart \\
Intent Classes  & \texttt{DATA\_QUERY} · \texttt{FOLLOWUP} · \texttt{SMALLTALK} \\
Self-correction & Tối đa 3 lần retry với error context \\
Post-insight    & Groq gọi thêm 1 lần để sinh 2--3 câu business insight \\
\bottomrule
\end{tabularx}
\end{tcolorbox}

\vfill
\begin{center}
\small\color{gray}
Olist Data Lakehouse Platform — Script Thuyết Trình Chi Tiết
\end{center}
\end{titlepage}

\tableofcontents
\newpage

% ============================================================
\section{Tổng quan — Agentic BI với Groq}
% ============================================================

\subsection{Vị trí trong hệ thống}

Agentic BI là \textbf{layer cuối cùng} trong Medallion pipeline — nhận Gold tables
(đã aggregate và partitioned) và cho phép người dùng hỏi bằng ngôn ngữ tự nhiên
thay vì viết SQL thủ công.

\begin{center}
\begin{tikzpicture}[
  node distance=0.5cm and 1.0cm,
  box/.style={draw, rounded corners=4pt, minimum width=2.4cm,
              minimum height=0.8cm, align=center, font=\small, line width=1pt},
  arr/.style={-Stealth, line width=1.5pt},
]
\node[box, fill=bronze!25, draw=bronze]              (b) {BRONZE\\{\tiny Iceberg}};
\node[box, fill=silver!40, draw=silver!80, right=of b] (s) {SILVER\\{\tiny PySpark}};
\node[box, fill=gold!40,   draw=gold!90,  right=of s] (g) {GOLD\\{\tiny 5 tables}};
\node[box, fill=groqgreen!20, draw=groqgreen, right=1.5cm of g,
      minimum width=3.0cm] (bi)
  {\textbf{Agentic BI}\\{\tiny Groq + DuckDB}};

\draw[arr, bronze]      (b) -- (s);
\draw[arr, silver!80]   (s) -- (g);
\draw[arr, groqgreen]   (g) -- node[font=\tiny, above]{NL query} (bi);

\node[font=\footnotesize, color=groqgreen, below=0.2cm of bi]
  {Streamlit UI · Chat interface};
\end{tikzpicture}
\end{center}

\subsection{Tại sao Groq?}

\begin{tcolorbox}[colback=groqgreen!5, colframe=groqgreen!50, arc=5pt]
\renewcommand{\arraystretch}{1.4}
\begin{tabularx}{\textwidth}{>{\bfseries}p{3.5cm} X}
\toprule
Điểm mạnh & Lý do chọn \\
\midrule
Tốc độ siêu nhanh &
  Groq LPU (Language Processing Unit) — không phải GPU, thiết kế chuyên biệt
  cho inference. Output token: $>$500 tok/s vs GPU $\sim$60 tok/s. \\
Free tier đủ dùng &
  6000 req/ngày, 500K token/phút — không tốn chi phí cho demo và prototype. \\
LLaMA 3.3-70B &
  Open-weight model Meta, hỗ trợ tiếng Việt tốt, giỏi sinh SQL,
  context 128K token. \\
OpenAI-compatible &
  API \texttt{/v1/chat/completions} — dùng chung \texttt{requests.post}, không cần SDK riêng. \\
\bottomrule
\end{tabularx}
\end{tcolorbox}

\newpage

% ============================================================
\section{Kiến trúc Pipeline — 4 Phase}
% ============================================================

\begin{tcolorbox}[colback=gray!5, colframe=groqgreen!50, arc=6pt,
  title={\faSitemap\ Sơ đồ — Groq Agentic BI: Text-based SQL Pipeline}]

\begin{center}
\begin{tikzpicture}[
  node distance=0.45cm and 0.85cm,
  box/.style={
    draw, rounded corners=5pt, minimum width=2.8cm,
    minimum height=0.85cm, align=center, font=\small, line width=1pt
  },
  phase/.style={
    draw, rounded corners=3pt, minimum width=1.0cm,
    minimum height=0.5cm, align=center, font=\tiny\bfseries,
    fill=white, line width=0.8pt
  },
  arr/.style={-Stealth, line width=1.3pt},
  darr/.style={-Stealth, line width=0.8pt, dashed},
]

% ── Row 1: main happy path ────────────────────────────────────
\node[box, fill=gray!15, draw=gray]               (user)
  {User\\{\tiny câu hỏi NL}};

\node[box, fill=intentblue!15, draw=intentblue,
      right=of user]                               (intent)
  {Intent Classifier\\{\tiny Groq — classify}};

\node[box, fill=sqlorange!20, draw=sqlorange,
      right=of intent]                             (llm)
  {LLaMA 3.3-70B\\{\tiny Groq API → JSON}};

\node[box, fill=duckdbgold!30, draw=duckdbgold,
      right=of llm]                                (duck)
  {DuckDB\\{\tiny execute SQL}};

\node[box, fill=insightgreen!20, draw=insightgreen,
      right=of duck]                               (insight)
  {Summarize\\{\tiny Groq 2nd call}};

\node[box, fill=groqgreen!15, draw=groqgreen,
      right=of insight]                            (render)
  {Streamlit\\{\tiny chart + insight}};

% ── Phase labels (above) ──────────────────────────────────────
\node[phase, draw=intentblue, text=intentblue,
      above=0.6cm of intent] (p1l) {Phase 1};
\node[phase, draw=sqlorange, text=sqlorange,
      above=0.6cm of llm]    (p2l) {Phase 2};
\node[phase, draw=duckdbgold, text=duckdbgold,
      above=0.6cm of duck]   (p3l) {Phase 3};
\node[phase, draw=insightgreen, text=insightgreen,
      above=0.6cm of insight] (p4l) {Phase 4};

\draw[dotted, intentblue]  (p1l) -- (intent);
\draw[dotted, sqlorange]   (p2l) -- (llm);
\draw[dotted, duckdbgold]  (p3l) -- (duck);
\draw[dotted, insightgreen](p4l) -- (insight);

% ── Happy path arrows ─────────────────────────────────────────
\draw[arr, gray]          (user)    -- (intent);
\draw[arr, intentblue]    (intent)  -- node[font=\tiny,above]{DATA\_QUERY} (llm);
\draw[arr, sqlorange]     (llm)     -- node[font=\tiny,above]{JSON} (duck);
\draw[arr, duckdbgold]    (duck)    -- node[font=\tiny,above]{df} (insight);
\draw[arr, insightgreen]  (insight) -- (render);

% ── Self-correction loop ──────────────────────────────────────
\node[box, fill=red!10, draw=red!50,
      below=1.2cm of duck]                        (retry)
  {Self-correct\\{\tiny retry với error msg}};

\draw[darr, red!60]  (duck.south)  -- node[font=\tiny,right]{ERROR} (retry.north);
\draw[darr, red!60]  (retry.west)  to[bend right=30]
                     node[font=\tiny,below]{fix SQL (max 3x)} (llm.south);

% ── Smalltalk / Followup shortcuts ───────────────────────────
\node[box, fill=llamapurple!15, draw=llamapurple,
      below=1.2cm of intent]                      (small)
  {SMALLTALK /\\FOLLOWUP};

\draw[darr, llamapurple]  (intent.south)  --
      node[font=\tiny,right]{other} (small.north);
\draw[darr, llamapurple]  (small.east)   to[bend left=20]  (render.south);

% ── Gold Iceberg ──────────────────────────────────────────────
\node[box, fill=gold!30, draw=gold!80,
      below=1.2cm of llm]                         (gold)
  {Gold Iceberg\\{\tiny 5 tables · PyIceberg}};

\draw[darr, gold!70]  (gold.east) to[bend right=20]
      node[font=\tiny,below]{schema ctx} (llm.south);
\draw[darr, gold!70]  (gold.east) to[bend left=10]  (duck.south);

% ── Retry counter badge ───────────────────────────────────────
\node[draw=red!40, fill=red!5, rounded corners=2pt,
      font=\tiny, right=0.3cm of retry]
  {attempt 1/2/3};

\end{tikzpicture}
\end{center}

\begin{center}
\small
\begin{tabular}{clll}
\toprule
Phase & Tên & Groq Call & Output \\
\midrule
1 & Intent Classify  & \texttt{classify\_intent()} — Groq     & DATA\_QUERY / FOLLOWUP / SMALLTALK \\
2 & SQL Generation   & \texttt{\_generate\_text\_based()} — Groq   & JSON \texttt{\{sql, explanation, chart\_type\}} \\
3 & Execute          & DuckDB in-memory (không gọi Groq)   & \texttt{pd.DataFrame} \\
4 & Summarize        & \texttt{\_summarize\_result()} — Groq       & 2--3 câu business insight \\
\bottomrule
\end{tabular}
\end{center}

\end{tcolorbox}

\newpage

% ============================================================
\section{Phase 1 — Intent Classification}
% ============================================================

\begin{phasebox}{intentblue}{\faTag\ Phase 1: Intent Classification}

Trước khi gọi SQL pipeline, hệ thống phân loại câu hỏi thành một trong ba loại:

\vspace{0.5em}
\begin{tabularx}{\textwidth}{>{\bfseries\ttfamily}p{3.2cm} X p{2.5cm}}
\toprule
Intent & Xử lý & Ví dụ \\
\midrule
DATA\_QUERY    & Chạy full SQL pipeline (Phase 2 → 4) &
  ``Doanh thu theo tháng?'' \\
FOLLOWUP       & Groq trả lời dựa trên lịch sử chat, không query DB &
  ``Sao bang SP lại cao vậy?'' \\
SMALLTALK      & Groq trả lời trực tiếp, không query DB &
  ``Bạn là AI gì?'' \\
\bottomrule
\end{tabularx}

\vspace{0.8em}

\textbf{Cách hoạt động:} \texttt{classify\_intent()} gọi Groq với system prompt
mô tả ba intent và rules. LLaMA 3.3-70B trả về một từ:
\texttt{DATA\_QUERY}, \texttt{FOLLOWUP}, hoặc \texttt{SMALLTALK}.
Nếu parse lỗi, mặc định fallback về \texttt{DATA\_QUERY}.

\end{phasebox}

% ============================================================
\section{Phase 2 — SQL Generation (Groq LLaMA 3.3-70B)}
% ============================================================

\begin{phasebox}{sqlorange}{\faBolt\ Phase 2: SQL Generation — Text-based JSON}

\subsection*{System Prompt Strategy}

Thay vì dùng native \textbf{Tool Use API} (như Anthropic), Groq path dùng
\textbf{text-based prompting}: yêu cầu LLM trả về \textit{thuần JSON},
sau đó parse bằng Python.

\vspace{0.5em}

\textbf{System prompt bao gồm:}
\begin{itemize}[leftmargin=1.5em, itemsep=3pt]
  \item Luật SQL: chỉ SELECT, tự thêm \texttt{LIMIT 1000}, tên bảng available
  \item Schema đầy đủ 5 Gold tables (tên cột, kiểu dữ liệu) — inject lúc runtime
  \item Quy tắc chart type: \texttt{bar/line/pie/scatter/none}
  \item \textbf{6 few-shot examples} (Q/A pair) — giúp LLM học pattern SQL chuẩn
  \item Luật đặc biệt cho sentiment table (\texttt{review\_sentiment}) — không dùng
        \texttt{review\_score} để suy ra sentiment
  \item Format output: JSON thuần, không markdown, không dấu backtick
\end{itemize}

\vspace{0.5em}

\textbf{Output format mong đợi:}
\begin{codebox}[title={JSON output từ LLaMA 3.3-70B}]
\begin{verbatim}
{
  "sql": "SELECT customer_state, COUNT(*) AS orders ...",
  "explanation": "Top 10 bang có nhiều đơn hàng nhất.",
  "chart_type": "bar"
}
\end{verbatim}
\end{codebox}

\subsection*{Few-shot Examples (trích từ code)}

\begin{codebox}[title={Ví dụ few-shot trong system prompt}]
\begin{verbatim}
Q: Doanh thu theo tháng năm 2018?
A: {"sql": "SELECT strftime(purchased_at, '%Y-%m') AS month,
    ROUND(SUM(payment_value),2) AS revenue_brl FROM fct_orders
    WHERE purchased_at >= '2018-01-01' AND purchased_at < '2019-01-01'
    GROUP BY 1 ORDER BY 1",
    "explanation": "Tổng doanh thu BRL theo tháng 2018.",
    "chart_type": "line"}

Q: Tỷ lệ các loại thanh toán?
A: {"sql": "SELECT payment_type, COUNT(*) AS cnt FROM fct_orders
    WHERE payment_type IS NOT NULL GROUP BY payment_type ORDER BY cnt DESC",
    "explanation": "Số đơn hàng theo phương thức thanh toán.",
    "chart_type": "pie"}
\end{verbatim}
\end{codebox}

\end{phasebox}

\newpage

% ============================================================
\section{Phase 3 — DuckDB Execution \& Self-correction}
% ============================================================

\begin{phasebox}{duckdbgold}{\faDatabase\ Phase 3: Execute + Self-correction Loop}

\subsection*{DuckDB In-memory Views}

DuckDB đọc dữ liệu trực tiếp từ Gold Iceberg tables thông qua PyIceberg:

\begin{enumerate}[leftmargin=1.5em, itemsep=4pt]
  \item PyIceberg load 5 Gold tables → Arrow format
  \item Đăng ký mỗi table thành \textbf{DuckDB in-memory view}
        (\texttt{con.register("fct\_orders", df)})
  \item SQL từ LLaMA chạy trực tiếp trên views — không cần Spark, không cần Trino
  \item Kết quả trả về \texttt{pd.DataFrame} trong vài milliseconds
\end{enumerate}

\vspace{0.8em}

\subsection*{Self-correction Loop (tối đa 3 lần)}

\begin{center}
\begin{tikzpicture}[
  node distance=0.4cm and 1.2cm,
  corrbox/.style={draw, rounded corners=4pt, minimum width=3.5cm,
               minimum height=0.75cm, align=center, font=\small, line width=1pt},
  arr/.style={-Stealth, line width=1.2pt},
]
\node[corrbox, fill=sqlorange!20, draw=sqlorange] (gen)
  {LLM sinh JSON \& SQL\\{\tiny attempt 1}};

\node[corrbox, fill=duckdbgold!25, draw=duckdbgold, right=of gen] (exec)
  {DuckDB execute};

\node[corrbox, fill=insightgreen!20, draw=insightgreen, right=of exec] (ok)
  {\faCheck\ Success\\{\tiny return DataFrame}};

\node[corrbox, fill=red!15, draw=red!50, below=1.0cm of exec] (fail)
  {SQL Error\\{\tiny parse error / column not found}};

\node[corrbox, fill=sqlorange!15, draw=sqlorange!70,
      left=of fail] (fix)
  {Retry Groq\\{\tiny SQL + error msg}};

\draw[arr, sqlorange]  (gen)  -- (exec);
\draw[arr, duckdbgold] (exec) -- node[font=\tiny,above]{OK} (ok);
\draw[arr, red!60]     (exec) -- node[font=\tiny,right]{ERROR} (fail);
\draw[arr, red!60]     (fail) -- (fix);
\draw[arr, sqlorange]  (fix)  to[bend left=30]
      node[font=\tiny,above]{retry} (exec);

\node[font=\tiny, color=red!60, below=0.2cm of fail]
  {Sau 3 lần: trả lỗi về UI};
\end{tikzpicture}
\end{center}

\vspace{0.5em}

\textbf{Fix prompt} khi SQL lỗi:
\begin{codebox}[title={Self-correction prompt gửi cho Groq}]
\begin{verbatim}
SQL sau gặp lỗi khi thực thi trên DuckDB:

SQL: SELECT seller_id, AVG(review_score) ...
Lỗi: Binder Error: "review_score" not found. Did you mean "avg_review_score"?

Hãy sửa lại SQL. Chỉ trả về JSON (không markdown):
{"sql": "...", "explanation": "...", "chart_type": "..."}
\end{verbatim}
\end{codebox}

\end{phasebox}

% ============================================================
\section{Phase 4 — Business Insight Summarization}
% ============================================================

\begin{phasebox}{insightgreen}{\faLightbulb\ Phase 4: Summarize Result}

Sau khi DuckDB trả về DataFrame, Groq được gọi \textbf{lần thứ hai} để sinh
2--3 câu business insight:

\begin{itemize}[leftmargin=1.5em, itemsep=4pt]
  \item Input: câu hỏi gốc + top 10 dòng đầu từ DataFrame
  \item Output: 2--3 câu tiếng Việt, nêu số liệu cụ thể với đơn vị đầy đủ
  \item Quy tắc: không dùng ``có thể'', ``dường như'', ``khoảng'' — chỉ fact cứng
  \item Tiền BRL: format ``X,XX tỷ R\$'' / ``X,XX triệu R\$''
  \item Highlight giá trị cao nhất / thấp nhất / bất thường
\end{itemize}

\vspace{0.5em}
\textbf{Ví dụ output:}
\begin{tcolorbox}[colback=insightgreen!8, colframe=insightgreen, arc=4pt]
\textit{``São Paulo dẫn đầu với 41.746 đơn hàng, chiếm 41,6\% tổng số đơn
toàn quốc. Rio de Janeiro đứng thứ 2 với 12.852 đơn --- cách biệt gần 3,3x
so với SP. Hai bang này cộng lại đã chiếm hơn 54\% thị phần đơn hàng Olist.''}
\end{tcolorbox}

\end{phasebox}

\newpage

% ============================================================
\section{Script Demo Chi Tiết}
% ============================================================

\begin{speakerbox}{p4color}{\faUser\ Người 4 · Phần Agentic BI (khoảng 2 phút)}

\subsection*{[Mở Streamlit Agentic BI] \timing{11:30 -- 13:30}}

\begin{demobox}
\textbf{Chuẩn bị:}
\begin{itemize}[leftmargin=1.5em, itemsep=2pt]
  \item Đảm bảo \texttt{AI\_PROVIDER=groq} và \texttt{GROQ\_API\_KEY} đã set trong \texttt{.env}
  \item Mở \texttt{http://localhost:8501} $\to$ chọn page \textbf{Olist Agentic BI}
  \item Sidebar sẽ hiển thị: \texttt{LLM: llama-3.3-70b-versatile} và \texttt{Mode: Text-based SQL}
  \item 5 Gold tables đều phải có dấu \textcolor{green!60!black}{\faCheckCircle} trong sidebar
\end{itemize}
\end{demobox}

\vspace{0.5em}

\subsubsection*{Demo 1 — Câu hỏi đơn: Top doanh thu theo bang \timing{11:30 -- 12:10}}

\begin{demobox}
\textbf{Gõ vào chat input:}
\begin{quote}
\textit{``Top 10 bang có doanh thu cao nhất năm 2018?''}
\end{quote}
\textbf{Quan sát:}
\begin{enumerate}[leftmargin=1.5em, itemsep=2pt]
  \item Spinner ``Đang xử lý...'' hiện ra $\to$ ``Phân tích câu hỏi và sinh SQL...''
  \item Sau $\sim$1 giây (Groq rất nhanh): tag \texttt{DATA QUERY} xanh lam xuất hiện
  \item Click expand \textbf{SQL} để show câu query được sinh ra
  \item Bar chart xuất hiện bên dưới
  \item Insight text: 2--3 câu về doanh thu SP, RJ, MG
\end{enumerate}
\end{demobox}

\begin{scriptbox}
``Đây là \textbf{Agentic BI} --- chat interface để hỏi dữ liệu Olist bằng tiếng Việt
thay vì viết SQL.

Backend dùng \textbf{Groq API} với model \textbf{LLaMA 3.3-70B} --- open-weight model
của Meta, chạy trên phần cứng LPU của Groq, nhanh hơn GPU thông thường từ 5--10 lần.

(expand SQL) Nhìn vào đây --- LLM tự sinh câu GROUP BY, tự filter năm 2018, tự ORDER BY
revenue, tự chọn bar chart vì đây là bài toán ranking. Nhóm em không hard-code gì ---
toàn bộ logic SQL đến từ model.

Kết quả: São Paulo dẫn đầu xa --- gấp hơn 3 lần Rio de Janeiro. Đây là insight
business thực tế từ 100 nghìn đơn hàng thật.''
\end{scriptbox}

\subsubsection*{Demo 2 — So sánh tỷ lệ giao trễ \timing{12:10 -- 12:50}}

\begin{demobox}
\textbf{Gõ tiếp vào chat:}
\begin{quote}
\textit{``So sánh tỷ lệ giao trễ giữa các phương thức thanh toán''}
\end{quote}
\textbf{Quan sát:}
\begin{enumerate}[leftmargin=1.5em, itemsep=2pt]
  \item SQL sinh ra sẽ có \texttt{CASE WHEN delivery\_status='late' THEN 1 ELSE 0 END}
  \item Chart type: bar (so sánh danh mục)
  \item Insight: phương thức nào có tỷ lệ trễ cao nhất?
\end{enumerate}
\end{demobox}

\begin{scriptbox}
``Câu hỏi lần này phức tạp hơn --- cần tính tỷ lệ phần trăm. LLM tự sinh
\texttt{CASE WHEN} để tính late orders, rồi chia cho total để ra tỷ lệ.

Đây là điểm quan trọng: người dùng không cần biết schema, không cần biết
\texttt{delivery\_status='late'} là cột gì hay bảng nào. LLM đọc schema context
và tự điền vào.

Nếu SQL sai --- ví dụ tên cột không đúng --- hệ thống \textbf{tự retry} tối đa
3 lần, mỗi lần gửi lại error message cho Groq để sửa. Người dùng không thấy
gì --- chỉ thấy kết quả cuối cùng.''
\end{scriptbox}

\subsubsection*{Demo 3 — Follow-up câu hỏi \timing{12:50 -- 13:30}}

\begin{demobox}
\textbf{Gõ tiếp (không cần context mới):}
\begin{quote}
\textit{``Phương thức nào đáng lo ngại nhất? Tại sao?''}
\end{quote}
\textbf{Quan sát:}
\begin{enumerate}[leftmargin=1.5em, itemsep=2pt]
  \item Tag \texttt{FOLLOW-UP} màu xanh lá xuất hiện (không phải DATA QUERY)
  \item \textbf{Không có SQL} --- Groq trả lời dựa trên kết quả câu trước trong history
  \item Câu trả lời phân tích narrative từ data đã có
\end{enumerate}
\end{demobox}

\begin{scriptbox}
``Câu này không có SQL mới --- hệ thống nhận diện đây là \textbf{FOLLOWUP} intent,
chỉ cần diễn giải data cũ.

Groq đọc lịch sử 10 tin nhắn gần nhất và trả lời như một analyst --- nêu cụ thể
phương thức nào và con số tỷ lệ, không bịa số mới.

Ba intent class này --- DATA\_QUERY, FOLLOWUP, SMALLTALK --- cho phép hệ thống
\textbf{không lãng phí} Groq call vào database khi không cần thiết,
đồng thời duy trì context conversation tự nhiên.''
\end{scriptbox}

\end{speakerbox}

\newpage

% ============================================================
\section{Sơ đồ Chi Tiết — Groq Text-based vs Claude Tool Use}
% ============================================================

\begin{tcolorbox}[colback=gray!5, colframe=gray!40, arc=6pt,
  title={\faSitemap\ So sánh hai chế độ: Groq (Text-based) vs Claude (Tool Use)}]

\begin{center}
\begin{tikzpicture}[
  node distance=0.4cm and 0.7cm,
  box/.style={draw, rounded corners=4pt, minimum width=2.6cm,
              minimum height=0.75cm, align=center, font=\footnotesize, line width=1pt},
  arr/.style={-Stealth, line width=1.1pt},
  darr/.style={-Stealth, line width=0.8pt, dashed},
  label/.style={font=\tiny, color=gray},
]

% ── Groq path (top row) ───────────────────────────────────────
\node[font=\small\bfseries\color{groqgreen}] (gLabel)
  at (0, 2.5) {Groq (AI\_PROVIDER=groq)};

\node[box, fill=groqgreen!15, draw=groqgreen]
  (g_user) at (0, 1.6) {User Question};
\node[box, fill=sqlorange!20, draw=sqlorange,
      right=of g_user]    (g_llm)  {Groq LLaMA\\{\tiny → JSON}};
\node[box, fill=duckdbgold!25, draw=duckdbgold,
      right=of g_llm]     (g_duck) {DuckDB\\{\tiny execute}};
\node[box, fill=insightgreen!20, draw=insightgreen,
      right=of g_duck]    (g_sum)  {Groq 2nd\\{\tiny summarize}};
\node[box, fill=groqgreen!15, draw=groqgreen,
      right=of g_sum]     (g_out)  {Result\\{\tiny chart+insight}};

\draw[arr, groqgreen]   (g_user) -- (g_llm);
\draw[arr, sqlorange]   (g_llm)  -- node[label,above]{\{sql,chart\}} (g_duck);
\draw[arr, duckdbgold]  (g_duck) -- node[label,above]{df} (g_sum);
\draw[arr, insightgreen](g_sum)  -- (g_out);

% Self-correction annotation
\node[font=\tiny, color=red!60, below=0.15cm of g_duck]
  {\faSync\ retry x3 on error};

% ── Claude path (bottom row) ──────────────────────────────────
\node[font=\small\bfseries\color{blue!70}] (cLabel)
  at (0, -0.3) {Claude (AI\_PROVIDER=anthropic)};

\node[box, fill=blue!10, draw=blue!60]
  (c_user) at (0, -1.1) {User Question};
\node[box, fill=purple!15, draw=purple!60,
      right=of c_user]   (c_claude) {Claude\\{\tiny Tool Use API}};
\node[box, fill=orange!20, draw=orange!60,
      right=of c_claude] (c_tool)   {query\_db\\{\tiny tool call}};
\node[box, fill=duckdbgold!25, draw=duckdbgold,
      right=of c_tool]   (c_duck)   {DuckDB};
\node[box, fill=blue!10, draw=blue!60,
      right=of c_duck]   (c_out)    {Result\\{\tiny + insight}};

\draw[arr, blue!60]    (c_user)   -- (c_claude);
\draw[arr, purple!60]  (c_claude) -- node[label,above]{tool\_use} (c_tool);
\draw[arr, orange!60]  (c_tool)   -- (c_duck);
\draw[arr, duckdbgold] (c_duck)   to[bend right=20]
      node[label,below]{tool\_result} (c_claude);
\draw[arr, blue!60]    (c_claude) -- (c_out);

\node[font=\tiny, color=purple!70, below=0.15cm of c_claude]
  {agentic loop: multi-call, self-corrects natively};

% ── Divider ───────────────────────────────────────────────────
\draw[gray!40, line width=1pt, dashed]
  (-1.5, 0.55) -- (12.5, 0.55);

\end{tikzpicture}
\end{center}

\vspace{0.5em}
\renewcommand{\arraystretch}{1.5}
\begin{tabularx}{\textwidth}{>{\bfseries}p{3.5cm}XX}
\toprule
 & \textcolor{groqgreen}{\bfseries Groq (Text-based)} & \textcolor{blue!70}{\bfseries Claude (Tool Use)} \\
\midrule
Số Groq/Claude call & 2 (generate + summarize) & 1+ (agentic loop) \\
Multi-query & Không (single-shot) & Có (gọi tool nhiều lần) \\
Self-correction & Explicit retry loop & Native (thấy ERROR trong tool result) \\
Chi phí & Free tier đủ dùng & $\sim$\$0.02/câu hỏi \\
Tốc độ đầu tiên & $<$1s (LPU) & 2--3s \\
Complex analysis & Hạn chế (1 query) & Mạnh (multi-step) \\
Code dependency & Chỉ \texttt{requests} & \texttt{anthropic} SDK \\
Fallback mode & Mặc định nếu không có \texttt{ANTHROPIC\_KEY} & Cần \texttt{ANTHROPIC\_API\_KEY} \\
\bottomrule
\end{tabularx}

\end{tcolorbox}

\newpage

% ============================================================
\section{Code Reference — Đường đi thực tế}
% ============================================================

\subsection{Config — Chọn provider}

\begin{codebox}[title={bi/config.py — Groq settings}]
\begin{verbatim}
# .env
AI_PROVIDER = groq
GROQ_API_KEY = gsk_xxxxxxxxxxxxxxxxxxxx
GROQ_MODEL   = llama-3.3-70b-versatile   # default

# config.py reads env and exposes:
AI_PROVIDER  = "groq"
GROQ_KEY     = "gsk_..."
GROQ_MODEL   = "llama-3.3-70b-versatile"
MAX_ROWS     = 1000
MAX_RETRIES  = 3
\end{verbatim}
\end{codebox}

\subsection{\_call\_llm — HTTP call tới Groq}

\begin{codebox}[title={bi/agent.py — Groq API call (OpenAI-compatible endpoint)}]
\begin{verbatim}
def _call_llm(messages, system, max_tokens=1024):
    # Groq path
    url   = "https://api.groq.com/openai/v1/chat/completions"
    model = GROQ_MODEL  # "llama-3.3-70b-versatile"
    key   = GROQ_KEY

    all_msgs = [{"role": "system", "content": system}] + messages
    resp = requests.post(
        url,
        headers={"Authorization": f"Bearer {key}"},
        json={"model": model, "messages": all_msgs,
              "max_tokens": max_tokens, "temperature": 0},
        timeout=60,
    )
    return resp.json()["choices"][0]["message"]["content"]
\end{verbatim}
\end{codebox}

\subsection{\_generate\_text\_based — SQL pipeline với retry}

\begin{codebox}[title={bi/agent.py — Text-based SQL + self-correction}]
\begin{verbatim}
def _generate_text_based(question, history, schema, con):
    system = _SYSTEM_SQL.format(
        max_rows=MAX_ROWS,
        table_names=", ".join(GOLD_TABLES),
        schema=schema,
    )
    base_msgs = history[-MAX_HISTORY:] + [{"role":"user","content":question}]

    for attempt in range(MAX_RETRIES):   # tối đa 3 lần
        if attempt == 0:
            msgs = base_msgs
        else:
            # Kèm SQL lỗi + error message → Groq tự fix
            msgs = base_msgs + [
                {"role": "assistant", "content": prev_raw},
                {"role": "user",      "content": FIX_MSG.format(
                    sql=raw_sql, error=last_error)}
            ]

        raw    = _call_llm(msgs, system)        # Groq call
        parsed = _parse_json(raw)               # Extract JSON
        sql    = parsed["sql"] + " LIMIT 1000"
        df     = con.execute(sql).df()          # DuckDB execute
        return {**parsed, "result_df": df}      # success!
        # on Exception: last_error = str(e), retry
\end{verbatim}
\end{codebox}

\subsection{Streamlit UI — Status và render}

\begin{codebox}[title={bi/agentic\_bi.py — Streamlit processing}]
\begin{verbatim}
with st.status("Đang xử lý...", expanded=True) as status:
    if AI_PROVIDER == "groq":
        status.write("Phân tích câu hỏi và sinh SQL...")

    res = run(intent=intent, question=user_input,
              history=chat_history, schema=schema_ctx, con=con)

    status.update(label="Hoàn thành", state="complete")

# Render: intent tag + SQL expander + chart + insight box
_render_assistant({...})
\end{verbatim}
\end{codebox}

\newpage

% ============================================================
\section{Sơ đồ Đầy Đủ — Groq Agentic BI End-to-End}
% ============================================================

\begin{tcolorbox}[colback=gray!5, colframe=groqgreen!50, arc=6pt,
  title={\faSitemap\ End-to-end: User → Groq → DuckDB → Gold Iceberg → Streamlit}]

\begin{center}
\begin{tikzpicture}[
  node distance=0.5cm and 0.7cm,
  box/.style={
    draw, rounded corners=5pt, minimum width=2.5cm,
    minimum height=0.85cm, align=center, font=\footnotesize, line width=1pt
  },
  arr/.style={-Stealth, line width=1.3pt},
  darr/.style={-Stealth, line width=0.8pt, dashed, gray!60},
]

% ── Row 1: Input ─────────────────────────────────────────────
\node[box, fill=gray!15, draw=gray] (browser)
  {Browser\\{\tiny Streamlit UI}};
\node[box, fill=intentblue!15, draw=intentblue,
      right=of browser] (validator)
  {Validator\\{\tiny length/injection}};
\node[box, fill=intentblue!20, draw=intentblue,
      right=of validator] (classifier)
  {Intent\\Classifier};

% ── Row 2 (DATA_QUERY path) ───────────────────────────────────
\node[box, fill=gray!10, draw=gray!50,
      below=0.8cm of classifier] (schema_ctx)
  {Schema\\Context\\{\tiny 5-table DDL}};
\node[box, fill=sqlorange!20, draw=sqlorange,
      right=of classifier] (groq1)
  {Groq\\{\tiny LLaMA call 1}\\{\tiny → JSON}};
\node[box, fill=yellow!20, draw=orange!50,
      right=of groq1] (parser)
  {JSON\\Parser\\{\tiny \_parse\_json()}};

% ── Row 3: DuckDB ─────────────────────────────────────────────
\node[box, fill=duckdbgold!30, draw=duckdbgold,
      right=of parser] (duckdb)
  {DuckDB\\{\tiny execute}};
\node[box, fill=red!10, draw=red!40,
      below=0.8cm of duckdb] (retry_box)
  {Self-correct\\{\tiny retry×3}};

% ── Row 4: Output ─────────────────────────────────────────────
\node[box, fill=insightgreen!20, draw=insightgreen,
      right=of duckdb] (groq2)
  {Groq\\{\tiny LLaMA call 2}\\{\tiny → insight}};
\node[box, fill=purple!10, draw=purple!50,
      right=of groq2] (charts)
  {Plotly\\Chart};
\node[box, fill=groqgreen!20, draw=groqgreen,
      right=of charts] (output)
  {Streamlit\\Output};

% ── Gold Iceberg (bottom) ─────────────────────────────────────
\node[box, fill=gold!30, draw=gold!80,
      below=1.5cm of groq1, xshift=1.5cm] (iceberg)
  {Gold Iceberg\\{\tiny 5 tables · R2}};
\node[box, fill=duckdbgold!20, draw=duckdbgold,
      right=of iceberg] (arrow_views)
  {DuckDB\\Views\\{\tiny PyIceberg→Arrow}};

% ── Arrows ────────────────────────────────────────────────────
\draw[arr, gray]          (browser)    -- (validator);
\draw[arr, intentblue]    (validator)  -- node[font=\tiny,above]{valid} (classifier);
\draw[arr, intentblue]    (classifier) -- node[font=\tiny,above]{DATA\_QUERY} (groq1);
\draw[darr]               (schema_ctx) -- (groq1);
\draw[arr, sqlorange]     (groq1)      -- (parser);
\draw[arr, orange!60]     (parser)     -- node[font=\tiny,above]{sql} (duckdb);
\draw[arr, duckdbgold]    (duckdb)     -- node[font=\tiny,above]{df} (groq2);
\draw[arr, insightgreen]  (groq2)      -- (charts);
\draw[arr, groqgreen]     (charts)     -- (output);

% Self-correction
\draw[darr, red!70] (duckdb.south) -- node[font=\tiny,right]{ERROR} (retry_box.north);
\draw[darr, red!70] (retry_box.west) to[bend right=35]
      node[font=\tiny,below]{fix prompt} (groq1.south);

% Iceberg
\draw[darr, gold!70]     (iceberg)    -- (arrow_views);
\draw[darr, duckdbgold]  (arrow_views) to[bend left=20] (duckdb.south);

% Schema context
\draw[darr] (iceberg.west) to[bend right=30]
      node[font=\tiny,left]{DDL} (schema_ctx.south);

% Followup shortcut
\node[box, fill=llamapurple!15, draw=llamapurple,
      below=0.8cm of classifier] (followup)
  {FOLLOWUP /\\SMALLTALK};
\draw[darr, llamapurple]  (classifier.south) -- (followup.north);
\draw[darr, llamapurple]  (followup.east) to[bend left=15] (groq1.south);

\end{tikzpicture}
\end{center}

\end{tcolorbox}

\newpage

% ============================================================
\section{Demo Questions — Cheat Sheet}
% ============================================================

\begin{tcolorbox}[colback=groqgreen!5, colframe=groqgreen!50, arc=6pt,
  title={\faList\ Câu hỏi demo + Expected Output}]

\renewcommand{\arraystretch}{1.6}
\begin{tabularx}{\textwidth}{>{\ttfamily\small}p{5.5cm} p{2.5cm} X}
\toprule
\normalfont\bfseries Câu hỏi & \bfseries Chart & \bfseries Expected Insight \\
\midrule
Top 10 bang doanh thu cao nhất 2018? &
  bar &
  SP $\gg$ RJ, chiếm $>$40\% \\

So sánh tỷ lệ giao trễ theo phương thức thanh toán &
  bar &
  boleto có tỷ lệ trễ cao hơn credit\_card \\

Doanh thu theo tháng năm 2018? &
  line &
  peak tháng 11 (Black Friday Brazil) \\

Phân tích churn: \% khách churned? &
  pie &
  $\sim$97\% churned (1-time buyers) \\

Tỷ lệ sentiment của review? &
  pie &
  positive $\approx$75\%, negative $\approx$13\% \\

Seller nào review tốt nhất $>$100 đơn? &
  bar &
  top sellers: avg score $>$4.8 \\

Tỷ lệ chuyển đổi lead theo kênh? &
  bar &
  organic\_search vs paid \\

So sánh Q1/2017 vs Q1/2018 &
  bar &
  2018 tăng mạnh, growth $>$100\% \\
\midrule
\multicolumn{3}{l}{\bfseries Follow-up (sau câu query trước):} \\
``Sao SP lại cao vậy?'' & none (followup) & Groq phân tích từ data đã có \\
``Xu hướng gì đáng lo?'' & none (followup) & Groq highlight anomaly \\
\bottomrule
\end{tabularx}

\end{tcolorbox}

\vspace{1em}

\begin{tcolorbox}[colback=yellow!8, colframe=orange!60, arc=4pt,
  title={\faExclamationTriangle\ Lưu ý khi demo}]
\begin{itemize}[leftmargin=1.5em, itemsep=4pt]
  \item \textbf{Groq free tier}: 30 req/min, 6000 req/day --- đủ cho demo 15 phút
  \item \textbf{Lần đầu khởi động}: PyIceberg load 5 tables vào DuckDB mất $\sim$10s
        (chỉ một lần duy nhất, sau đó cached)
  \item \textbf{Nếu câu hỏi quá phức tạp}: Groq text-based chỉ sinh 1 SQL.
        Demo nên dùng câu hỏi đơn rõ ràng, tránh ``phân tích toàn diện...''
  \item \textbf{Sidebar check}: trước khi demo verify 5 bảng \textcolor{green!60!black}{\faCheckCircle}:
        \texttt{fct\_orders}, \texttt{fct\_funnel}, \texttt{dim\_sellers},
        \texttt{dim\_customers}, \texttt{review\_sentiment}
  \item \textbf{Nếu SQL lỗi}: hệ thống tự retry, spinner hiển thị. Chờ thêm 2--3s.
  \item \textbf{Mode label}: sidebar hiển thị ``Mode: Text-based SQL'' --- đây là Groq mode.
        Nếu thấy ``Tool Use'' nghĩa là đang dùng Claude mode.
\end{itemize}
\end{tcolorbox}

% ============================================================
\section{Phụ lục — Q\&A về Agentic BI}
% ============================================================

\begin{tcolorbox}[colback=blue!5, colframe=blue!40, arc=6pt,
  title={\faQuestion\ Câu hỏi thường gặp — Agentic BI}]

\renewcommand{\arraystretch}{1.6}
\begin{tabularx}{\textwidth}{>{\bfseries}p{5cm} X}
\toprule
Câu hỏi & Trả lời \\
\midrule
Sao không dùng Claude Tool Use? &
  Groq \textbf{miễn phí} và \textbf{nhanh hơn} cho single-query use case.
  Claude Tool Use mạnh hơn ở multi-step analysis nhưng tốn tiền.
  Hệ thống \textbf{hỗ trợ cả hai} --- chỉ cần đổi \texttt{AI\_PROVIDER} trong \texttt{.env}. \\

Tại sao không dùng text-to-SQL riêng (vd: Langchain)? &
  Không cần thêm dependency. Groq với system prompt + few-shot đủ mạnh,
  self-correction tự xử lý 95\% SQL lỗi. \\

Có thể hỏi join nhiều bảng không? &
  Có. Schema context gồm cả 5 bảng. LLM có thể sinh JOIN.
  Ví dụ: ``Seller state nào có review sentiment tốt nhất?''
  → JOIN \texttt{dim\_sellers} + \texttt{review\_sentiment}. \\

DuckDB có chậm không khi data lớn? &
  DuckDB in-memory với 100K rows Gold tables: $<$100ms.
  Bottleneck là Groq latency ($\sim$800ms) --- nhưng vẫn nhanh hơn nhiều so với
  Spark hay Trino cho ad-hoc query nhỏ. \\

SQL injection risk? &
  System prompt enforce ``chỉ SELECT''.
  Ngoài ra validator ở \texttt{bi/validator.py} chặn các pattern
  nguy hiểm trước khi gửi cho LLM. \\

Groq có hỗ trợ tiếng Việt tốt không? &
  LLaMA 3.3-70B hỗ trợ tiếng Việt tốt ở mức production.
  System prompt và few-shot examples đều tiếng Việt.
  Insight output bằng tiếng Việt, số liệu đúng đơn vị. \\
\bottomrule
\end{tabularx}

\end{tcolorbox}

% ============================================================
\end{document}
% ============================================================
